% ============================================================================
%  Presentación (slides) — Climate Cognitive System (CCS)
%  Curso: Computación Cognitiva — UTEC
%  Compilar:  tectonic slides.tex        (recomendado, sin dependencias)
%         o:  latexmk -pdf slides.tex    /    pdflatex slides.tex (x2)
% ============================================================================
\documentclass[aspectratio=169,11pt]{beamer}

% Engine-agnostic font setup: works with pdflatex AND xelatex/lualatex/tectonic.
\usepackage{iftex}
\ifPDFTeX
  \usepackage[utf8]{inputenc}
  \usepackage[T1]{fontenc}
  \usepackage{lmodern}
  \usepackage{textcomp}
\else
  \usepackage{fontspec}
\fi
\usepackage[spanish]{babel}
\usepackage{booktabs}
\usepackage{tikz}
\usetikzlibrary{shapes.geometric,arrows.meta,positioning,calc}
\usepackage{hyperref}

% ---- Theme & branding ------------------------------------------------------
\usetheme{Boadilla}
\definecolor{ccsblue}{HTML}{1F6FEB}
\definecolor{ccsdark}{HTML}{0B2545}
\definecolor{ccssoft}{HTML}{E8F0FE}
\usecolortheme[named=ccsblue]{structure}
\setbeamercolor{title}{fg=ccsdark}
\setbeamercolor{frametitle}{fg=white,bg=ccsdark}
\setbeamertemplate{navigation symbols}{}
\setbeamertemplate{itemize items}[circle]

% ---- Metadata --------------------------------------------------------------
\title[Climate Cognitive System]{Climate Cognitive System}
\subtitle{Gestión climática predictiva para aulas universitarias}
\author[Grupo CCS]{Balcazar Santa Cruz, Kiara \and Gonzales Torres, Aaron \and Surco Vergara, Maria Fernanda}
\institute[UTEC]{Computación Cognitiva --- UTEC, Barranco (Lima)\\ \footnotesize Docente: Farfán Madariaga, Jaime Moshe}
\date{2 de julio de 2026}

\begin{document}

% ============================ TÍTULO =======================================
\begin{frame}[plain]
  \titlepage
\end{frame}

% ============================ PROBLEMA =====================================
\begin{frame}{El problema}
  Los aires acondicionados (A/C) de las aulas operan en modo \textbf{always-on} por horario.
  \begin{itemize}
    \item Enfrían aulas \textbf{vacías o con baja asistencia} con la misma intensidad.
    \item No consideran la \textbf{ocupación real} ni las condiciones ambientales.
    \item HVAC $\approx$ \textbf{50\,\%} del consumo eléctrico de un edificio
      no residencial (Pérez-Lombard et al., 2008).
    \item Resultado: \textbf{sobreconsumo energético} sistemático y sin trazabilidad.
  \end{itemize}
  \vfill
  \begin{block}{¿A quién afecta?}
    A la institución (sobrecosto), al mantenimiento (sin datos) y al medio ambiente
    (huella de carbono).
  \end{block}
\end{frame}

% ============================ SOLUCIÓN =====================================
\begin{frame}{La propuesta: CCS}
  Plataforma IoT de climatización \textbf{predictiva} que ajusta el \emph{setpoint}
  del A/C combinando ocupación real, clima y un motor cognitivo (ML + heurístico).
  \begin{columns}[T]
    \begin{column}{0.5\textwidth}
      \textbf{Control tradicional}
      \begin{itemize}
        \item Always-on por horario
        \item Reactivo
        \item Ignora la ocupación
      \end{itemize}
    \end{column}
    \begin{column}{0.5\textwidth}
      \textbf{CCS (feed-forward + feedback)}
      \begin{itemize}
        \item Pre-enfría con la reserva
        \item Corrige con CO\textsubscript{2} en vivo
        \item Calcula ROI energético
      \end{itemize}
    \end{column}
  \end{columns}
\end{frame}

% ============================ PERFILES =====================================
\begin{frame}{Perfiles del sistema (RBAC)}
  \begin{itemize}
    \item \textbf{Usuario} (\texttt{guest} / \texttt{collaborator}): inicia sesión,
      consulta aulas/sensores en tiempo real, crea reservas, recibe alertas de CO;
      el \texttt{collaborator} controla el sensor \emph{con reserva activa}.
    \item \textbf{Administrador} (\texttt{admin}): gestiona usuarios, aprovisiona
      aulas y sensores, supervisa telemetría, revisa ROI, recarga el modelo de IA.
  \end{itemize}
  \vfill
  \footnotesize Cuentas \texttt{guest} auto-aprobadas; \texttt{collaborator}/\texttt{admin} quedan \texttt{pending}.
\end{frame}

% ============================ APP + HARDWARE ===============================
\begin{frame}{Aplicación web y hardware}
  \begin{columns}[T]
    \begin{column}{0.5\textwidth}
      \textbf{App web (React + Vite)}
      \begin{itemize}\small
        \item Login / registro (JWT)
        \item Dashboard en vivo
        \item Aulas, reservas, infraestructura
        \item Reporte de ROI + chat Watson
      \end{itemize}
    \end{column}
    \begin{column}{0.5\textwidth}
      \textbf{Hardware (Arduino UNO)}
      \begin{itemize}\small
        \item DHT11 --- temperatura / humedad
        \item MQ-135 --- CO\textsubscript{2} $\rightarrow$ ocupación
        \item MQ-7 --- monóxido de carbono (ADC $\rightarrow$ ppm)
        \item Cámara + MediaPipe --- aforo
        \item Emisor IR --- actuador del A/C (fase 2)
      \end{itemize}
      \footnotesize Serial $\rightarrow$ \texttt{hardware/sistema.py} $\rightarrow$ \texttt{POST /sensors/} cada ${\sim}2$\,s.
    \end{column}
  \end{columns}
\end{frame}

% ============================ IA ===========================================
\begin{frame}{Inteligencia Artificial (núcleo, no adorno)}
  \begin{itemize}
    \item \textbf{Predicción de carga térmica:} \texttt{HistGradientBoostingRegressor}
      sobre ocupación esperada/real, clima y metadata del aula.
    \item \textbf{Ocupación real desde CO\textsubscript{2}:} balance de masa
      (\texttt{personas} $\approx$ (CO\textsubscript{2} $-$ base) / ppm-persona).
    \item \textbf{Fallback heurístico}, \textbf{detección de anomalías} (CO),
      \textbf{visión por computadora} (aforo por cámara + MediaPipe),
      \textbf{asistente NLP} (IBM Watson) y \textbf{reportes de ROI}.
  \end{itemize}
  \vfill
  \begin{block}{Rigor}
    Holdout cronológico + baselines (media/heurística): el modelo se despliega
    \textbf{sólo si supera} a la heurística (\texttt{beats\_baselines}).
  \end{block}
\end{frame}

% ============================ ARQUITECTURA =================================
\begin{frame}{Arquitectura de la solución}
  \begin{center}
  \begin{tikzpicture}[
    font=\scriptsize, node distance=6mm and 10mm,
    box/.style={draw=ccsdark,fill=ccssoft,rounded corners,minimum height=7mm,minimum width=18mm,align=center},
    actor/.style={draw=ccsblue,fill=white,rounded corners,minimum height=7mm,minimum width=16mm,align=center},
    ia/.style={draw=ccsblue,fill=ccsblue!15,rounded corners,minimum height=7mm,minimum width=18mm,align=center},
    hw/.style={draw=gray,fill=white,rounded corners,minimum height=7mm,minimum width=20mm,align=center},
    arr/.style={-{Stealth[length=2mm]},gray}]
    \node[actor](u){Usuario};
    \node[actor,below=of u](a){Admin};
    \node[box,right=of $(u)!0.5!(a)$](app){App Web\\React};
    \node[box,right=of app](api){Backend\\FastAPI};
    \node[box,above right=4mm and 10mm of api](mongo){MongoDB};
    \node[box,below right=4mm and 10mm of api](pg){PostgreSQL};
    \node[ia,right=28mm of api](ml){Modelo ML\\+ Watson};
    \node[hw,below=12mm of app](hw){Arduino UNO + Cámara};
    \draw[arr](u)--(app); \draw[arr](a)--(app);
    \draw[arr]([yshift=1mm]app.east)--([yshift=1mm]api.west);
    \draw[arr]([yshift=-1mm]api.west)--([yshift=-1mm]app.east);
    \draw[arr](api)--(mongo); \draw[arr](api)--(pg); \draw[arr](api)--(ml);
    \draw[arr](hw.east)-|(api.south) node[midway,below,black]{\tiny POST /sensors/};
  \end{tikzpicture}
  \end{center}
\end{frame}

% ============================ TECNOLOGÍAS ==================================
\begin{frame}{Tecnologías}
  \small
  \begin{tabular}{@{}ll@{}}
    \toprule
    \textbf{Capa} & \textbf{Stack} \\
    \midrule
    Frontend & React 18, Vite, TypeScript, Tailwind, Recharts \\
    Backend  & FastAPI, SQLAlchemy; JWT + Argon2 \\
    Datos    & PostgreSQL 15, MongoDB 6 \\
    IA       & scikit-learn, IBM Watson Assistant, MediaPipe \\
    Hardware & Arduino UNO, DHT11, MQ-135, MQ-7, cámara \\
    Deploy   & Docker + Docker Compose, Nginx \\
    \bottomrule
  \end{tabular}
\end{frame}

% ============================ VIABILIDAD + IMPACTO =========================
\begin{frame}{Viabilidad e impacto}
  \begin{columns}[T]
    \begin{column}{0.5\textwidth}
      \textbf{Viabilidad}
      \begin{itemize}\small
        \item HW de bajo costo; un nodo funcional.
        \item Riesgo: pocos sensores.
        \item Mitigación: simulador + \textbf{dataset sembrado reproducible}.
      \end{itemize}
    \end{column}
    \begin{column}{0.5\textwidth}
      \textbf{Impacto}
      \begin{itemize}\small
        \item Ambiental: menos kWh y CO\textsubscript{2}.
        \item Económico: ahorro en factura.
        \item Institucional: trazabilidad del uso real de aulas.
        \item Seguridad: alertas de CO.
        \item Automatización basada en datos.
      \end{itemize}
    \end{column}
  \end{columns}
\end{frame}

% ============================ CIERRE / ENLACES ============================
\begin{frame}[plain]{Recursos del proyecto}
  \vfill
  \centering
  {\Large\textbf{\color{ccsdark}¡Gracias!}}\\[1.2cm]

  \begin{tabular}{rl}
    \textbf{Repositorio GitHub:} &
      \href{https://github.com/Kiarosaurus/Climate-Cognitive-System}{\color{ccsblue}github.com/Kiarosaurus/Climate-Cognitive-System} \\[6pt]
    \textbf{Instancia desplegada:} &
      \href{https://kiarosaurus.me}{\color{ccsblue}kiarosaurus.me} \\[6pt]
    \textbf{Videos (SW y HW):} &
      \href{https://drive.google.com/drive/folders/1n6j4DU4aDTeXRAdlSHInQun2IR3O-uuh?usp=sharing}{\color{ccsblue}Google Drive --- Demos de Software y Hardware} \\
  \end{tabular}
  \vfill
\end{frame}

\end{document}
